%%%%%%%%%%%%%%%%%%%%%%%%%%%%%%%%%%%%%%%%%
% Twenty Seconds Resume/CV
% LaTeX Template
% Version 1.1 (8/1/17)
%
% This template has been downloaded from:
% http://www.LaTeXTemplates.com
%
% Original author:
% Carmine Spagnuolo (cspagnuolo@unisa.it) with major modifications by
% Vel (vel@LaTeXTemplates.com)
%
% License:
% The MIT License (see included LICENSE file)
%
%%%%%%%%%%%%%%%%%%%%%%%%%%%%%%%%%%%%%%%%%

% ----------------------------------------------------------------------------------------
%	PACKAGES AND OTHER DOCUMENT CONFIGURATIONS
% ----------------------------------------------------------------------------------------

\documentclass[letterpaper]{../cls/twentysecondcv} % a4paper for A4
\usepackage{hyperref}

% Definitions of Pinyin
\makeatletter
\def\py@yunpriv#1{%
  \if a#1 10\else
    \if o#1 9\else
      \if e#1 8\else
        \if i#1 7\else
          \if u#1 6\else
            \if v#1 5\else
              \if A#1 4\else
                \if O#1 3\else
                  \if E#1 2\fi\fi\fi\fi\fi\fi\fi\fi\fi0
  }

  \def\py@init{%
    \edef\py@befirst{}%
    \edef\py@char{}\edef\py@tuneletter{}%
    \def\py@last{}%
    \def\py@tune{5}%
  }

  % Usage:
  % \pinyin{Hao3hao3\ xue2xi2} （好好学习）
  \def\pinyin#1{%
    \edef\py@postscan{#1}%
    \py@init
    % scan
    \loop
    \edef\py@char{\expandafter\@car\py@postscan\@nil}%
    \edef\py@postscan{\expandafter\@cdr\py@postscan\@nil}%
    \ifnum 0 < 0\py@char
      \edef\py@tune{\py@char}%
      \py@first \py@tuneat\py@tuneletter\py@tune \py@last\kern -4sp\kern 4sp{}\py@init
    \else
      \ifnum\py@yunpriv\py@char > \py@yunpriv\py@tuneletter
        \edef\py@tuneletter{\py@char}\edef\py@first{\py@befirst}\def\py@last{}%
      \else
        \edef\py@last{\py@last\if v\py@char\"u\else\py@char\fi}%
      \fi
      \edef\py@befirst{\py@befirst\if v\py@char\"u\else\py@char\fi}%
    \fi
    \ifx\py@postscan\@empty\else
      \repeat
    }

    \let\py@macron \=
    \let\py@acute \'
    \let\py@hacek \v
    \let\py@grave \`

    %% \py@tuneat{Letter}{tune}
    \def\py@tuneat#1#2{%
      \if v#1%
        \py@tune@v #2%
      \else
        \if i#1%
          \py@tune@i #2%
        \else
          \ifcase#2%
          \or\py@macron #1\or\py@acute #1\or\py@hacek #1\or\py@grave #1\else #1%
          \fi
        \fi\fi
    }

    \def\py@tune@v#1{{%
        \dimen@ii 1ex%
        \fontdimen5\font 1.1ex%
        \rlap{\"u}%
        \fontdimen5\font .6ex%
        \ifcase#1%
        \or\py@macron u\or\py@acute u\or\py@hacek u\or\py@grave u\else u%
        \fi
        \fontdimen5\font\dimen@ii
      }}

    \def\py@tune@i#1{%
      \ifcase#1
      \or\py@macron \i\or\py@acute \i\or\py@hacek \i\or\py@grave \i\else i%
      \fi
    }
    \makeatletter
    % End of pinyin

    % ----------------------------------------------------------------------------------------
    % PERSONAL INFORMATION
    % ----------------------------------------------------------------------------------------

    % If you don't need one or more of the below, just remove the content leaving the command, e.g. \cvnumberphone{}

    \profilepic{../img/W1-2} % Profile picture
    \virtcv{../img/virtcv}
    \cvname{Pedro Gomes \\Branquinho} % Your name
    \cvjobtitle{Engenheiro Físico} % Job title/career
    \cvdate{07 Outubro 1997} % Date of birth
    \cvaddress{Franca, SP} % Short address/location, use \newline if more than 1 line is required
    \cvnumberphone{(16) 99340-1215} % Phone number
    \linkedin{\href{https://www.linkedin.com/in/pedro-g-branquinho/}{Perfil LinkedIn}} % Personal website
    \cvsite{\href{www.buddhilw.com}{www.buddhilw.com}} % Personal website
    \cvmail{pedro.branquinho@usp.br} % Email address

    % ----------------------------------------------------------------------------------------

    \begin{document}

    % ----------------------------------------------------------------------------------------
    % ABOUT ME
    % ----------------------------------------------------------------------------------------

    \aboutme{Sou um Desenvolvedor de Software, com formação em Engenharia. Tenho conhecimento em ambos Front e Back- end, bem como conhecimento avançado de Estatística e Computação Científica.

      Tenho experiência free-lance na Indústria, em empresas de pequeno e grande porte. Atualmente, trabalho com as tecnologias baseadas em Go, Python, Julia, Java e JavaScript (FullStack)}

    % To have no About Me section, just remove all the text and leave \aboutme{}

    % ----------------------------------------------------------------------------------------
    % SKILLS
    % ----------------------------------------------------------------------------------------
    \skills{{Python/5.5}, {Go/5},{Julia/4.5},{Git/4.5},{Clojure(Script)/4.5},
      {Scripting (Bash, et al.)/4}, {JavaScript/4.5}}

    % ----------------------------------------------------------------------------------------

    \makeprofile % Print the sidebar

    % ----------------------------------------------------------------------------------------
    % INTERESTS
    % ----------------------------------------------------------------------------------------

    \section{Interesses Profissionais}

    % \vspace{10mm}

    \begin{itemize}
      \item Full-stack Development.
      \item Microservices.
      \item DevOps.
      \item Data Science.
      % \item Data Science - Teste de Hipótese, Sistemas de Recomendação, PCA;
      % \item Aplicações Web, usando Machine Learning;
      % \item Automação de pipelines Linux/Unix;
      % \item Desenvolvimento Backend;
    \end{itemize}
    \vspace{0.2cm}
    % ----------------------------------------------------------------------------------------
    % EDUCATION
    % ----------------------------------------------------------------------------------------

    \section{Educação}

    \begin{twenty} % Environment for a list with descriptions
      \twentyitem{2016 - 2022}{Graduação (Completa)}{\normalsize{Universidade de São Paulo, USP}}{Engenharia Física}
    \end{twenty}
    \vspace{0.2cm}
    % ----------------------------------------------------------------------------------------
    % PUBLICATIONS
    % ----------------------------------------------------------------------------------------

    \section{Publicações Eletônicas}

    \begin{twentyshort} % Environment for a short list with no
      % descriptions
      \twentyitemshort{2021}{\href{https://github.com/BuddhiLW/commons-csv-clj}{\color{blue!50!violet}
          Indústria (Flow Finance) - Automação de Boletos}}
      \twentyitemshort{2021}{\href{https://github.com/BuddhiLW/CloshBashika}{\color{blue!50!violet}
          Industria (Lupo S.A.) - Automação de Relatórios Técnicos com Clojure/LaTeX.}}
      % \twentyitemshort{2020}{\href{https://github.com/BuddhiLW/EEL_EstatMultiVarRN}{\color{blue!50!violet}Machine
      % Learning, com R - Enquanto cursava Estatística Multivareada.}}
      % \twentyitemshort{2020}{\href{https://www.youtube.com/watch?v=spYCKElN-v0&list=PLweBgessvalf9eGUknfHQZ6_q2Nz7t8df}{\color{blue!50!violet}
      % Um Minicurso em \LaTeX{}} - Disponível no YouTube.}
    \end{twentyshort}
    \vspace{0.2cm}

    % ----------------------------------------------------------------------------------------
    % Industry Experience
    % ----------------------------------------------------------------------------------------

    \vspace{5mm}
    \section{Jobs -- Industry and Academia}

    \begin{twenty} % Environment for a short list with no descriptions
      % \twentyitem{<dates>}{<title>}{<location>}{<description>}
      \twentyitem{2021}{Lupo S.A. -- Jan/2021-Maio/2021}{Desenvolvedor}{
        Escrevi um software para facilitar a escrita de relatórios técnicos, enquanto trabalhando na WJB Engenharia, com contrato com a Lupo S.A.. O trabalho foi feito, com Clojure.
      }
      \twentyitem{2021}{FlowFinance S.C.- Jun/2021-Julio/2021}{Developer}{
        Escrevi uma aplicação para gerar uma lista de pagamentos, para lançamentos de boletos. A empresa prestava esse serviço para BIORC. O trabalho foi feito, com Clojure.
      }
      \twentyitem{2021/2022}{University of São Paulo - Sep/2021-Mar/2022}{Researcher}{
        Pesquisei maneiras implementar soluções numéricas para modelos matemáticos de Tráfego de Carros (PDEs), com Julia e Python, usando I.A..
      }
      \twentyitem{2021/2022}{Café do Bem (NPO) - Agosto/2021-Atualmente}{Voluntariado}{Eu criei o site, \href{https://cafe-do-bem.company.site/}{https://cafe-do-bem.company.site/}, o qual é uma plataforma para venda de café. A ONG possui a missão de reverter os ganhos com populações de risco.
      }
      \twentyitem{2023}{FACTI -- Jan/2023-Agosto/2023}{FullStack Dev.}{
        Experiência desenvolvendo aplicação com Vanilla JavaScript, Node.js, Express.js, ClojureScript e um framework criado para resolver o problema da empresa, usando uma Linguagem criada de Domínio Específico (\textit{DSL}).
      }
    \end{twenty}

    \vspace{5mm}
    % ----------------------------------------------------------------------------------------
    % LANGUAGE
    % ----------------------------------------------------------------------------------------
    \vspace{0.2cm}
    \section{Línguas e Fluência}

    \begin{description}
      \item[Português:] Nativo.
      \item[Inglês:] \href{https://www.efset.org/cert/hqg62J}{Fluência a nível C2} (clique para abrir o certificado).
      \item[Mandarim:] HSK2.
    \end{description}

    \vspace{0.2cm}


    % ----------------------------------------------------------------------------------------
    % SECOND PAGE EXAMPLE
    % ----------------------------------------------------------------------------------------

    \newpage % Start a new page

    \makeprofileNoExtra % Print the sidebar

    \section{\LARGE{Certificados}}
    \subsection{\textbf{Google Go Specialization}}
    \begin{description}
      \item[\href{https://buddhilw.github.io/bug-free-fiesta/}{Getting Started with Go:}] Sintaxe da Língua e Estruturas de Dados.
      \item[\href{https://buddhilw.github.io/bug-free-fiesta/}{Functions, Methods and Interfaces:}] Orientação a Objeto, com Go.
      \item[\href{https://buddhilw.github.io/bug-free-fiesta/}{Concurrency in Go:}] Tópicos avançados em Go.
    \end{description}
    \vspace{10mm}
    \subsection{\textbf{Backend by Meta (former Facebook)}}
    \begin{description}
      \item[\href{https://buddhilw.github.io/bug-free-fiesta/}{Introduction to Back-End Development:}] Agile, TDD etc.; Infraestrutura da Web, HTTP, CTP e outros protocólogos.
      \item[\href{https://buddhilw.github.io/bug-free-fiesta/}{Version Control:}] Praticas com Git e Github.
      \item[\href{https://buddhilw.github.io/bug-free-fiesta/}{Programming in Python:}] Sintaxe e Orientação a Objetos.
    \end{description}
    \vspace{10mm}
    \subsection{\textbf{Responsive Web Design by freeCodeCamp}}

    \begin{description}
      \item[\href{https://buddhilw.github.io/bug-free-fiesta/}{Responsive Web Design:}] CSS; Bootstrap e outros Frameworks.
    \end{description}

    \vspace{10mm}
    \subsection{\textbf{Introduction to Docker: Build Your Own Portfolio Site by Coursera}}

    \begin{description}
      \item[\href{https://buddhilw.github.io/bug-free-fiesta/}{Introduction to Docker: Build Your Own Portfolio Site:}] Projeto prático criando aplicações em containers e lançando um website, com Docker.
    \end{description}

    \vspace{10mm}
    \subsection{\textbf{English certification by EF Education First}}

    \begin{description}
      \item[\href{https://buddhilw.github.io/bug-free-fiesta/}{EF SET Score - C2}] Teste de Inglês, com base no European Standards - 78 pontos, com nota C2, máximo intervalo.
      % English test, based on European Standards - 78 points, giving the C2 score, at the maximum-range.
    \end{description}
    \vspace{5mm}
    \section{Site agregador, com todas minhas Certificações}

    \vspace{1mm}

    \begin{center}
      Link: \href{https://buddhilw.github.io/bug-free-fiesta/}{https://buddhilw.github.io/bug-free-fiesta/}
      \vspace{1mm}

      \includegraphics[width=\imagewidth]{../img/virtcerts.png}
    \end{center}
    \newpage
    \makeprofileNoExtra

    \section{Publicações Eletrónicas}

    \begin{twentyshort} % Environment for a short list with no
      % descriptions
      \twentyitemshort{2020}{\href{https://github.com/BuddhiLW/EEL_EstatMultiVarRN}{\color{blue!50!violet}Machine
          Learning, with R - while enrolled in Multivariate Statics.}}
      \twentyitemshort{2020}{\href{https://github.com/BuddhiLW/MC-LaTeX/tree/master/LabEELw/LabEELw/MaterialMC/Livro_PT1}{\color{blue!50!violet}Introduction
          to \LaTeX, P1.}}
      \twentyitemshort{2020}{\href{https://github.com/BuddhiLW/MC-LaTeX/tree/master/LabEELw/LabEELw/MaterialMC/Livro_PT2}{\color{blue!50!violet}
          Essential Commands in \LaTeX, P2.
        }}
      \twentyitemshort{2020}{\href{https://github.com/BuddhiLW/MC-LaTeX/tree/master/LabEELw/LabEELw/MaterialMC/Livro_PT3}{\color{blue!50!violet}Bibliographical
          References, Citations and Beamer Presentations, P3.}}
      \twentyitemshort{2020}{\href{https://www.youtube.com/watch?v=spYCKElN-v0&list=PLweBgessvalf9eGUknfHQZ6_q2Nz7t8df}{\color{blue!50!violet}
          A Minicourse on \LaTeX{}} - Material available on YouTube (Portuguese).}
    \end{twentyshort}

    % ----------------------------------------------------------------------------------------
    % EXPERIENCE
    % ----------------------------------------------------------------------------------------

    \section{Experiencias Extracurriculares - Enquanto Universitário}

    \begin{twenty} % Environment for a list with descriptions
      \twentyitem{2022}{\textit{Programming in Python}}{Meta}{Curso sobre Algorítmos e Estrtura de dados.}
      \twentyitem{2022}{\textit{Version Control}}{Meta}{Curso em Git, Github e melhor práticas, em Versionamento.}
      \twentyitem{2022}{\textit{Introduction to Back-End Development}}{Meta}{Curso de introdução às melhores práticas, TDD, Git e Metodologias Ágeis.}
      \twentyitem{2022}{\textit{Google Go Specialization}}{University of California (UCI)}{Profissionalização em Go/Golang: Programação Orientada a Objetos (POO), Paralelismo, Concorrência e Design de Softwares.}
      \twentyitem{2021}{\textit{Marco: 50 projetos pessoas e 26 colaborações (PR)}}{GitHub}{Meu atual
        {\href{https://github.com/BuddhiLW/}{\color{blue!50!violet}status no GitHub}} (08/2021)}
      \twentyitem{2020}{\textit{International Congress on Funcional
          Programming (ICFP))}}{Penn University}{Maior congresso em linguagens funcionais - exposição ao estado da arte da computação.}
      \twentyitem{2020}{\textit{Introduction to Git and
          GitHub}}{Google}{Parte da especialização, \textit{Google IT Automation with Python}}
      \twentyitem{2019}{\textit{Clojure for the Brave and True}}{Livro Texto}{Início do auto-aprendizado de Clojure.}
      \twentyitem{2018}{Instalação do Arch Linux}{Arquitetura Linux}{Aprendizado de como configurar um Desktop, do zero. Experiência com APIs, integrações e Third-Party-Software.}
      \twentyitem{2017}{Emacs, SLIME, Common Lisp}{Open Source}{Quando comecei a estudar progamação, algoritmos funcionais e design de softwares.}
      \twentyitem{2017}{Ubuntu Linux}{Primeira Instalação}{Primeira experiência com Linux, SystemD e Shell Scripting.}
    \end{twenty}


  \end{document}
  %%% Local Variables:
  %%% TeX-master: t
  %%% End:
